\chapter{Implementasi Smart Contract Sistem DecMed}
\label{lampiran:implementasi-smart-contract-sistem-decmed}

Tabel~\ref{tab:daftar-module-smart-contract-decmed} sampai Tabel~\ref{tab:daftar-method-decmed-proxy} merinci modul Move dan method terkait yang digunakan pada implementasi \textit{smart contract} DecMed.

\let\textunderscorewithoutbreak\_
\newcommand{\breakabletextunderscore}{\textunderscorewithoutbreak\allowbreak{}}
\newcolumntype{Y}[1]{>{\let\_\breakabletextunderscore\raggedright\arraybackslash}p{#1}<{\let\_\textunderscorewithoutbreak}}

\begin{longtable}{|Y{0.36\textwidth}|Y{0.56\textwidth}|}
	\caption{Daftar Module Smart Contract Sistem DecMed}
	\label{tab:daftar-module-smart-contract-decmed}                                                                                                                                                    \\
	\hline
	\textbf{Nama Module}                 & \textbf{Deskripsi}                                                                                                                                          \\
	\hline
	\endfirsthead
	\hline
	\textbf{Nama Module}                 & \textbf{Deskripsi}                                                                                                                                          \\
	\hline
	\endhead

	\texttt{decmed::admin}               & Modul global admin untuk registrasi rumah sakit, aktivasi admin fasyankes, dan daftar rumah sakit.                                                          \\ \hline
	\texttt{decmed::hospital\_personnel} & Modul akun tenaga fasyankes, aktivasi, akses baca/tulis, \textit{delegation}, dan \textit{audit delegation}.                                                \\ \hline
	\texttt{decmed::patient}             & Modul akun pasien, \textit{grant} akses ke \texttt{AdministrativePersonnel}, log akses, \textit{audit delegation}, dan pembacaan metadata RME milik pasien. \\ \hline
	\texttt{decmed::proxy}               & Modul yang dipakai \textit{server} PRE untuk baca/tulis metadata RME.                                                                                       \\ \hline
	\texttt{decmed::shared}              & Modul penyedia \texttt{GlobalAdminCap} dan \texttt{ProxyCap}                                                                                                \\ \hline
\end{longtable}

Tabel~\ref{tab:daftar-module-smart-contract-decmed} menegaskan pemisahan tanggung jawab antarmodul untuk administrasi fasyankes, akun personel, pasien, proxy PRE, dan kapabilitas bersama.

\begin{longtable}{|Y{0.45\textwidth}|Y{0.47\textwidth}|}
	\caption{Daftar Method pada Module \texttt{decmed::admin}}
	\label{tab:daftar-method-decmed-admin}                                                                          \\
	\hline
	\textbf{Nama Method}             & \textbf{Deskripsi}                                                           \\
	\hline
	\endfirsthead
	\hline
	\textbf{Nama Method}             & \textbf{Deskripsi}                                                           \\
	\hline
	\endhead

	\texttt{create\_activation\_key} & Membuat data rumah sakit dan activation key pertama untuk admin rumah sakit. \\ \hline
	\texttt{update\_activation\_key} & Memperbarui activation key dan metadata admin rumah sakit.                   \\ \hline
	\texttt{get\_hospitals}          & Mengambil daftar rumah sakit untuk client kementerian/global admin.          \\ \hline
\end{longtable}

Tabel~\ref{tab:daftar-method-decmed-admin} merangkum operasi registrasi dan aktivasi fasyankes pada sisi admin global.

\begin{longtable}{|Y{0.45\textwidth}|Y{0.47\textwidth}|}
	\caption{Daftar Method pada Module \texttt{decmed::hospital\_personnel}}
	\label{tab:daftar-method-decmed-hospital-personnel}                                                                                \\
	\hline
	\textbf{Nama Method}                      & \textbf{Deskripsi}                                                                     \\
	\hline
	\endfirsthead
	\hline
	\textbf{Nama Method}                      & \textbf{Deskripsi}                                                                     \\
	\hline
	\endhead

	\texttt{create\_activation\_key}          & Admin fasyankes membuat activation key untuk tenaga administratif atau medis.          \\ \hline
	\texttt{use\_activation\_key}             & Menandai activation key tenaga fasyankes sebagai sudah digunakan.                      \\ \hline
	\texttt{signup}                           & Registrasi akun tenaga fasyankes ke address IOTA dan menyimpan metadata administratif. \\ \hline
	\texttt{is\_account\_registered}          & Mengecek apakah akun tenaga fasyankes sudah aktif/terdaftar.                           \\ \hline
	\texttt{get\_account\_state}              & Menentukan status \textit{onboarding} tenaga fasyankes.                                \\ \hline
	\texttt{get\_account\_info}               & Mengambil metadata akun, \textit{role}, \textit{subrole}, dan metadata rumah sakit.    \\ \hline
	\texttt{update\_account\_activation\_key} & Admin fasyankes mengganti activation key personel.                                     \\ \hline
	\texttt{update\_administrative\_metadata} & Memperbarui profil administratif tenaga fasyankes.                                     \\ \hline
	\texttt{get\_hospital\_personnels}        & Mengambil daftar personel di bawah admin fasyankes.                                    \\ \hline
	\texttt{get\_delegatee\_candidates}       & Mengambil kandidat personel yang bisa menerima delegasi akses.                         \\ \hline
	\texttt{cleanup\_read\_access}            & Membersihkan akses baca yang sudah kedaluwarsa.                                        \\ \hline
	\texttt{cleanup\_update\_access}          & Membersihkan akses \textit{update} yang sudah kedaluwarsa.                             \\ \hline
	\texttt{get\_read\_access}                & Mengambil daftar \textit{capability} baca milik tenaga fasyankes.                      \\ \hline
	\texttt{get\_update\_access}              & Mengambil daftar \textit{capability} \texttt{update} milik tenaga fasyankes.           \\ \hline
	\texttt{create\_delegated\_access}        & Membuat \textit{delegated access} untuk personel lain.                                 \\ \hline
	% \texttt{revoke\_delegated\_access}        & Mencabut delegated access dari sisi delegator tenaga fasyankes.                        \\ \hline
	\texttt{create\_medical\_record}          & Membuat entri RME blob besar. (\textit{legacy})                                        \\ \hline
\end{longtable}

Tabel~\ref{tab:daftar-method-decmed-hospital-personnel} menunjukkan cakupan method untuk aktivasi, pengelolaan akun, serta pembentukan dan pembersihan akses tenaga fasyankes.

\begin{longtable}{|Y{0.45\textwidth}|Y{0.47\textwidth}|}
	\caption{Daftar Method pada Module \texttt{decmed::patient}}
	\label{tab:daftar-method-decmed-patient}                                                                                   \\
	\hline
	\textbf{Nama Method}                      & \textbf{Deskripsi}                                                             \\
	\hline
	\endfirsthead
	\hline
	\textbf{Nama Method}                      & \textbf{Deskripsi}                                                             \\
	\hline
	\endhead

	\texttt{signup}                           & Registrasi akun pasien dan metadata administratif terenkripsi.                 \\ \hline
	\texttt{is\_account\_registered}          & Mengecek apakah \textit{address} pasien  terdaftar.                            \\ \hline
	\texttt{get\_account\_state}              & Menentukan status \textit{onboarding} pasien.                                  \\ \hline
	\texttt{get\_account\_info}               & Mengambil metadata administratif pasien.                                       \\ \hline
	\texttt{update\_administrative\_metadata} & Memperbarui metadata administratif pasien.                                     \\ \hline
	\texttt{get\_hospital\_personnel\_info}   & Mengambil info publik tenaga fasyankes saat pasien memproses QR.               \\ \hline
	\texttt{create\_access}                   & Memberi akses baca/tulis kepada tenaga fasyankes.                              \\ \hline
	\texttt{get\_access\_log}                 & Mengambil log \textit{grant} akses pasien.                                     \\ \hline
	\texttt{get\_delegation\_audit\_log}      & Mengambil log \textit{audit trail} \textit{delegation} pasien.                 \\ \hline
	\texttt{get\_medical\_records}            & Mengambil daftar metadata RME pasien.                                          \\ \hline
	\texttt{get\_medical\_record}             & Mengambil satu metadata RME berdasarkan indeks.                                \\ \hline
	\texttt{revoke\_access}                   & Mencabut akses tenaga fasyankes dan memicu \textit{cascade revoke delegation}. \\ \hline
\end{longtable}

Tabel~\ref{tab:daftar-method-decmed-patient} menegaskan peran pasien dalam pemberian akses, audit, dan pencabutan akses terhadap RME miliknya.

\begin{longtable}{|Y{0.45\textwidth}|Y{0.47\textwidth}|}
	\caption{Daftar Method pada Module \texttt{decmed::proxy}}
	\label{tab:daftar-method-decmed-proxy}                                                                                                          \\
	\hline
	\textbf{Nama Method}                          & \textbf{Deskripsi}                                                                              \\
	\hline
	\endfirsthead
	\hline
	\textbf{Nama Method}                          & \textbf{Deskripsi}                                                                              \\
	\hline
	\endhead

	\texttt{create\_capability}                   & Membuat \texttt{ProxyCap} untuk address service proxy.                                          \\ \hline
	\texttt{create\_medical\_record\_segment}     & Menambahkan metadata RME tersegmentasi ke daftar medical metadata pasien.                       \\ \hline
	\texttt{get\_administrative\_data}            & Mengambil metadata administratif pasien untuk kebutuhan akses/dekripsi.                         \\ \hline
	\texttt{get\_hospital\_personnel\_role}       & Mengecek role tenaga fasyankes saat setup token/key access.                                     \\ \hline
	\texttt{get\_hospital\_personnel\_auth\_info} & Mengecek \textit{role} dan \textit{subrole}.                                                    \\ \hline
	\texttt{get\_medical\_records}                & Mengambil daftar metadata RME pasien dan menyaringnya sebagai segmen berdasarkan cakupan akses. \\ \hline
	\texttt{get\_medical\_record}                 & Mengambil satu metadata segmen dan data pendukung untuk pembacaan \textit{payload}.             \\ \hline
	\texttt{is\_patient\_registered}              & Memvalidasi bahwa pasien terdaftar                                                              \\ \hline
	\texttt{create\_medical\_record}              & Membuat entri RME blob besar. (\textit{legacy})                                                 \\ \hline
	\texttt{update\_medical\_record}              & Melakukan edit entri RME blob besar. (\textit{legacy})                                          \\ \hline
	\texttt{get\_medical\_record\_update}         & Mengambil data RME untuk proses edit entri blob besar. (\textit{legacy})                        \\ \hline
	\texttt{get\_delegation\_access\_snapshot}    & Mengambil \textit{snapshot} metadata akses delegasi.                                            \\ \hline
\end{longtable}

Tabel~\ref{tab:daftar-method-decmed-proxy} menegaskan bahwa modul proxy menyediakan antarmuka \textit{on-chain} bagi \textit{Server} PRE untuk membaca dan menulis metadata RME sesuai otorisasi.
